% GENERATED FILE. Edit canonical lessons, then run npm run generate:books.
% canonical-source-hash: fnv1a64:aee07b1ea6cc99a2
% canonical-lessons: JA-C198-noru, JA-C198-oriru, JA-C198-jitensha, JA-C198-chikatetsu, JA-C198-tsuku

\chapter{Getting On and Off the Train}
\label{ch:ja-getting-on-and-off-the-train}
\hlchaptermodality{198}

\begin{chapteropening}
\textbf{By the end of this chapter:} \emph{I can say how I ride trains, bicycles and the subway.}

Complete the last lesson of chapter 198: I can say how I ride trains, bicycles and the subway.
\end{chapteropening}

\section[\texorpdfstring{\ja{のる} (noru)}{のる (noru)}]{\ja{のる} (noru) — to ride; to get on (a train or bus)}

\label{lesson:JA-C198-noru}



\begin{quote}
Before the new one: say the Japanese for open for business, then the Japanese for not yet open.
\end{quote}

\subsection*{You'll want to know: \ja{のる}}
\textbf{\ja{のる}} — \emph{noru} — \textquotedblleft{}to ride; to get on (a train or bus)\textquotedblright{}.

\begin{quote}
\textbf{\ja{まいあさ} \ja{でんしゃに} \ja{のります。}} — \emph{maiasa densha ni norimasu} — \textquotedblleft{}I take the train every morning.\textquotedblright{}
\end{quote}

Polite \textbf{\ja{のります}}, \emph{norimasu}; a godan verb.

\begin{cousinweb}
\textbf{\ja{のる}} (\emph{noru}) is a native Japanese word (wago), inherited rather than borrowed from Chinese; its kanji was attached later as a kun'yomi, and its older origin is not settled (JMdict, KANJIDIC2).
\end{cousinweb}

\begin{grammarlens}[title={What you've built}]
One more word for everyday life at A2.
\end{grammarlens}

\subsection*{Your turn}
\begin{itemize}
\raggedright
  \item \emph{Say it:} \emph{noru}
  \item \emph{Say it:} \emph{noru}, once more
  \item \emph{Say it:} \emph{maiasa densha ni norimasu} — I take the train every morning.
  \item \emph{Recall:} say the Japanese for open for business, then the Japanese for not yet open, then say \emph{noru} again
\end{itemize}

\begin{tcolorbox}[breakable,colback=teal!4,colframe=teal!35!black,title={Before you move on}]
What does \ja{のる} mean? (\textquotedblleft{}To ride; to get on (a train or bus)\textquotedblright{}.) Say it once more.
\end{tcolorbox}
\section[\texorpdfstring{\ja{おりる} (oriru)}{おりる (oriru)}]{\ja{おりる} (oriru) — to get off; to go down}

\label{lesson:JA-C198-oriru}



\begin{quote}
Before the new one: say the Japanese for not yet open, then the Japanese for to ride.
\end{quote}

\subsection*{You'll want to know: \ja{おりる}}
\textbf{\ja{おりる}} — \emph{oriru} — \textquotedblleft{}to get off; to go down\textquotedblright{}.

\begin{quote}
\textbf{\ja{ここで} \ja{でんしゃを} \ja{おります。}} — \emph{koko de densha o orimasu} — \textquotedblleft{}I get off the train here.\textquotedblright{}
\end{quote}

Polite \textbf{\ja{おります}}, \emph{orimasu}; an ichidan verb.

\begin{cousinweb}
\textbf{\ja{おりる}} (\emph{oriru}) is a native Japanese word (wago), inherited rather than borrowed from Chinese; its kanji was attached later as a kun'yomi, and its older origin is not settled (JMdict, KANJIDIC2).
\end{cousinweb}

\begin{grammarlens}[title={What you've built}]
One more word for everyday life at A2.
\end{grammarlens}

\subsection*{Your turn}
\begin{itemize}
\raggedright
  \item \emph{Say it:} \emph{oriru}
  \item \emph{Say it:} \emph{oriru}, once more
  \item \emph{Say it:} \emph{koko de densha o orimasu} — I get off the train here.
  \item \emph{Recall:} say the Japanese for not yet open, then the Japanese for to ride, then say \emph{oriru} again
\end{itemize}

\begin{tcolorbox}[breakable,colback=teal!4,colframe=teal!35!black,title={Before you move on}]
What does \ja{おりる} mean? (\textquotedblleft{}To get off; to go down\textquotedblright{}.) Say it once more.
\end{tcolorbox}
\section[\texorpdfstring{\ja{じてんしゃ} (jitensha)}{じてんしゃ (jitensha)}]{\ja{じてんしゃ} (jitensha) — a bicycle}

\label{lesson:JA-C198-jitensha}



\begin{quote}
Before the new one: say the Japanese for to ride, then the Japanese for to get off.
\end{quote}

\subsection*{You'll want to know: \ja{じてんしゃ}}
\textbf{\ja{じてんしゃ}} — \emph{jitensha} — \textquotedblleft{}a bicycle\textquotedblright{}.

\begin{quote}
\textbf{\ja{じてんしゃで} \ja{えきまで} \ja{いきます。}} — \emph{jitensha de eki made ikimasu} — \textquotedblleft{}I go to the station by bicycle.\textquotedblright{}
\end{quote}

\begin{cousinweb}
\textbf{\ja{じてんしゃ}} (\emph{jitensha}) is Sino-Japanese (kango): three Chinese characters read with their on'yomi, readings borrowed from Chinese. They are \emph{ji}, \textquotedblleft{}self\textquotedblright{} (Middle Chinese \emph{dzijH}), \emph{ten}, \textquotedblleft{}to turn\textquotedblright{} (\emph{trjwenX}) and \emph{sha}, \textquotedblleft{}cart\textquotedblright{} (\emph{tsyhae}) (KANJIDIC2; Middle Chinese in Baxter's transcription).
\end{cousinweb}

\begin{grammarlens}[title={What you've built}]
One more word for everyday life at A2.
\end{grammarlens}

\subsection*{Your turn}
\begin{itemize}
\raggedright
  \item \emph{Say it:} \emph{jitensha}
  \item \emph{Say it:} \emph{jitensha}, once more
  \item \emph{Say it:} \emph{jitensha de eki made ikimasu} — I go to the station by bicycle.
  \item \emph{Recall:} say the Japanese for to ride, then the Japanese for to get off, then say \emph{jitensha} again
\end{itemize}

\begin{tcolorbox}[breakable,colback=teal!4,colframe=teal!35!black,title={Before you move on}]
What does \ja{じてんしゃ} mean? (\textquotedblleft{}A bicycle\textquotedblright{}.) Say it once more.
\end{tcolorbox}
\section[\texorpdfstring{\ja{ちかてつ} (chikatetsu)}{ちかてつ (chikatetsu)}]{\ja{ちかてつ} (chikatetsu) — a subway; an underground railway}

\label{lesson:JA-C198-chikatetsu}



\begin{quote}
Before the new one: say the Japanese for to get off, then the Japanese for a bicycle.
\end{quote}

\subsection*{You'll want to know: \ja{ちかてつ}}
\textbf{\ja{ちかてつ}} — \emph{chikatetsu} — \textquotedblleft{}a subway; an underground railway\textquotedblright{}.

\begin{quote}
\textbf{\ja{ちかてつで} \ja{まちに} \ja{いきます。}} — \emph{chikatetsu de machi ni ikimasu} — \textquotedblleft{}I go into town on the subway.\textquotedblright{}
\end{quote}

\begin{cousinweb}
\textbf{\ja{ちかてつ}} (\emph{chikatetsu}) is Sino-Japanese (kango): three Chinese characters read with their on'yomi, readings borrowed from Chinese. They are \emph{chi}, \textquotedblleft{}earth\textquotedblright{} (Middle Chinese \emph{dijH}), \emph{ka}, \textquotedblleft{}below\textquotedblright{} (\emph{haeX}) and \emph{tetsu}, \textquotedblleft{}iron\textquotedblright{} (\emph{thet}) (KANJIDIC2; Middle Chinese in Baxter's transcription).
\end{cousinweb}

\begin{grammarlens}[title={What you've built}]
One more word for everyday life at A2.
\end{grammarlens}

\subsection*{Your turn}
\begin{itemize}
\raggedright
  \item \emph{Say it:} \emph{chikatetsu}
  \item \emph{Say it:} \emph{chikatetsu}, once more
  \item \emph{Say it:} \emph{chikatetsu de machi ni ikimasu} — I go into town on the subway.
  \item \emph{Recall:} say the Japanese for to get off, then the Japanese for a bicycle, then say \emph{chikatetsu} again
\end{itemize}

\begin{tcolorbox}[breakable,colback=teal!4,colframe=teal!35!black,title={Before you move on}]
What does \ja{ちかてつ} mean? (\textquotedblleft{}A subway; an underground railway\textquotedblright{}.) Say it once more.
\end{tcolorbox}
\section[\texorpdfstring{\ja{つく} (tsuku)}{つく (tsuku)}]{\ja{つく} (tsuku) — to get to; to arrive at}

\label{lesson:JA-C198-tsuku}



\begin{quote}
Before the new one: say the Japanese for a bicycle, then the Japanese for a subway.
\end{quote}

\subsection*{You'll want to know: \ja{つく}}
\textbf{\ja{つく}} — \emph{tsuku} — \textquotedblleft{}to get to; to arrive at\textquotedblright{}.

\begin{quote}
\textbf{\ja{でんしゃは} \ja{ゆうがた} \ja{えきに} \ja{つきます。}} — \emph{densha wa yūgata eki ni tsukimasu} — \textquotedblleft{}The train gets to the station in the evening.\textquotedblright{}
\end{quote}

Polite \textbf{\ja{つきます}}, \emph{tsukimasu}; a godan verb.

\begin{cousinweb}
\textbf{\ja{つく}} (\emph{tsuku}) is a native Japanese word (wago), inherited rather than borrowed from Chinese; its kanji was attached later as a kun'yomi, and its older origin is not settled (JMdict, KANJIDIC2).
\end{cousinweb}

\begin{grammarlens}[title={What you've built}]
One more word for everyday life at A2.
\end{grammarlens}

\subsection*{Your turn}
\begin{itemize}
\raggedright
  \item \emph{Say it:} \emph{tsuku}
  \item \emph{Say it:} \emph{tsuku}, once more
  \item \emph{Say it:} \emph{densha wa yūgata eki ni tsukimasu} — The train gets to the station in the evening.
  \item \emph{Recall:} say the Japanese for a bicycle, then the Japanese for a subway, then say \emph{tsuku} again
\end{itemize}

\begin{tcolorbox}[breakable,colback=teal!4,colframe=teal!35!black,title={Before you move on}]
What does \ja{つく} mean? (\textquotedblleft{}To get to; to arrive at\textquotedblright{}.) Say it once more.
\end{tcolorbox}
